
\documentclass[12pt]{article}
%%%%%%%%%%%%%%%%%%%%%%%%%%%%%%%%%%%%%%%%%%%%%%%%%%%%%%%%%%%%%%%%%%%%%%%%%%%%%%%%%%%%%%%%%%%%%%%%%%%%%%%%%%%%%%%%%%%%%%%%%%%%%%%%%%%%%%%%%%%%%%%%%%%%%%%%%%%%%%%%%%%%%%%%%%%%%%%%%%%%%%%%%%%%%%%%%%%%%%%%%%%%%%%%%%%%%%%%%%%%%%%%%%%%%%%%%%%%%%%%%%%%%%%%%%%%
\usepackage{amsmath}
\usepackage{amsfonts}
\usepackage{mitpress}

\setcounter{MaxMatrixCols}{10}
%TCIDATA{OutputFilter=LATEX.DLL}
%TCIDATA{Version=5.50.0.2953}
%TCIDATA{<META NAME="SaveForMode" CONTENT="1">}
%TCIDATA{BibliographyScheme=Manual}
%TCIDATA{Created=Friday, December 02, 2011 10:43:51}
%TCIDATA{LastRevised=Wednesday, February 15, 2012 14:00:22}
%TCIDATA{<META NAME="GraphicsSave" CONTENT="32">}
%TCIDATA{<META NAME="DocumentShell" CONTENT="Articles\SW\A Simple MIT Press Article">}
%TCIDATA{CSTFile=LaTeX Article (bright).cst}

\newtheorem{theorem}{Theorem}
\newtheorem{acknowledgement}[theorem]{Acknowledgement}
\newtheorem{algorithm}[theorem]{Algorithm}
\newtheorem{axiom}[theorem]{Axiom}
\newtheorem{case}[theorem]{Case}
\newtheorem{claim}[theorem]{Claim}
\newtheorem{conclusion}[theorem]{Conclusion}
\newtheorem{condition}[theorem]{Condition}
\newtheorem{conjecture}[theorem]{Conjecture}
\newtheorem{corollary}[theorem]{Corollary}
\newtheorem{criterion}[theorem]{Criterion}
\newtheorem{definition}[theorem]{Definition}
\newtheorem{example}[theorem]{Example}
\newtheorem{exercise}[theorem]{Exercise}
\newtheorem{lemma}[theorem]{Lemma}
\newtheorem{notation}[theorem]{Notation}
\newtheorem{problem}[theorem]{Problem}
\newtheorem{proposition}[theorem]{Proposition}
\newtheorem{remark}[theorem]{Remark}
\newtheorem{solution}[theorem]{Solution}
\newtheorem{summary}[theorem]{Summary}
\newenvironment{proof}[1][Proof]{\noindent\textbf{#1.} }{\ \rule{0.5em}{0.5em}}
\newdimen\dummy
\dummy=\oddsidemargin
\addtolength{\dummy}{72pt}
\marginparwidth=.5\dummy
\marginparsep=.1\dummy
\input{tcilatex}
\begin{document}

\title{The Title of A Simple MIT Press Style}
\author{A. U. Thor \\
%EndAName
At this Address}
\maketitle

\begin{abstract}
Replace this text with the text of your own abstract.
\end{abstract}

\section{The One Particle Idempotency Condition}

is given by%
\begin{align*}
& \mathbf{0}=\Gamma \left( \hat{\lambda}_{1}\right) ^{2}-\Gamma \left( \hat{%
\lambda}_{1}\right) =\left( \sum_{1\leq N\leq r}\Gamma _{N}\left( \hat{%
\lambda}_{1}^{2}\right) \right) ^{2}-\Gamma \left( \hat{\lambda}_{1}\right)
\\
& =\left( \sum_{1\leq N\leq r}N\Gamma _{1}^{N}\left( \hat{\lambda}%
_{1}\right) \right) ^{2}-\sum_{1\leq N\leq r}N\Gamma _{1}^{N}\left( \hat{%
\lambda}_{1}\right) \\
& =\sum_{1\leq N\leq r}N\Gamma _{1}^{N}\left( \hat{\lambda}_{1}\right)
\left( N\Gamma _{1}^{N}\left( \hat{\lambda}_{1}\right) -P_{\mathcal{H}%
^{N}}\right) \\
& =\quad \Gamma _{1}^{1}\left( \hat{\lambda}_{1}\right) ^{2}+4\Gamma
_{1}^{2}\left( \hat{\lambda}_{1}\right) ^{2}+9\Gamma _{1}^{3}\left( \hat{%
\lambda}_{1}\right) ^{2}+\cdots +r^{2}\Gamma _{1}^{r}\left( \hat{\lambda}%
_{1}\right) ^{2} \\
& -\Gamma _{1}^{1}\left( \hat{\lambda}_{1}\right) +2\Gamma _{1}^{2}\left( 
\hat{\lambda}_{1}\right) +3\Gamma _{1}^{3}\left( \hat{\lambda}_{1}\right)
+\cdots +r\Gamma _{1}^{r}\left( \hat{\lambda}_{1}\right)
\end{align*}%
Thus 
\begin{equation*}
\left\{ N\Gamma _{1}^{N}\left( \hat{\lambda}_{1}\right) \right\}
^{2}=N\Gamma _{1}^{N}\left( \hat{\lambda}_{1}\right) ;\,1\leq N\leq r 
\end{equation*}

\subsection{One Particle Condition}

The one particle condition is 
\begin{equation*}
\hat{\lambda}_{1}^{2}-\hat{\lambda}_{1}=0\Rightarrow \hat{\lambda}_{1jj}=1%
\text{ or }0;\,1\leq j\leq r. 
\end{equation*}

\subsection{Two Particle Condition}

The two particle expansion gives 
\begin{align*}
& \Gamma _{1}^{2}\left( \hat{\lambda}_{1}\right) =\hat{\lambda}_{1}\wedge P_{%
\mathcal{H}^{1}}=\sum_{\substack{ 1\leq j_{1}\leq 4  \\ 1\leq j_{2}\leq 4}}%
\lambda _{1j_{1}j_{1}}\left\vert \varphi _{j_{1}}\right\rangle \left\langle
\varphi _{j_{1}}\right\vert \wedge \left\vert \varphi _{j_{2}}\right\rangle
\left\langle \varphi _{j_{2}}\right\vert \\
& =\frac{1}{2}\sum_{1\leq j_{1}<j_{2}\leq 4}\left( \lambda
_{1j_{1}j_{1}}+\lambda _{1j_{2}j_{2}}\right) \left\vert \varphi
_{j_{1}}\varphi _{j_{2}}\right\rangle \left\langle \varphi _{j_{1}}\varphi
_{j_{2}}\right\vert
\end{align*}%
and its square is 
\begin{align*}
& \Gamma _{1}^{2}\left( \hat{\lambda}_{1}\right) ^{2}=\Gamma _{1}^{2}\left( 
\hat{\lambda}_{1}\right) \circ \Gamma _{1}^{2}\left( \hat{\lambda}%
_{1}\right) =\left( \hat{\lambda}_{1}\wedge P_{\mathcal{H}^{1}}\right) \circ
\left( \hat{\lambda}_{1}\wedge P_{\mathcal{H}^{1}}\right) \\
& =\frac{1}{4}\sum_{1\leq j_{1}<j_{2}\leq 4}\left\{ \left( \lambda
_{1j_{1}j_{1}}^{2}+\lambda _{1j_{2}j_{2}}^{2}\right) \right. +\left.
2\lambda _{1j_{1}j_{1}}\lambda _{1j_{2}j_{2}}\right\} \left\vert \varphi
_{j_{1}}\varphi _{j_{2}}\right\rangle \left\langle \varphi _{j_{1}}\varphi
_{j_{2}}\right\vert \\
& =\frac{1}{4}\sum_{1\leq j_{1}<j_{2}\leq 4}\left( \lambda
_{1j_{1}j_{1}}+\lambda _{1j_{2}j_{2}}\right) ^{2}\left\vert \varphi
_{j_{1}}\varphi _{j_{2}}\right\rangle \left\langle \varphi _{j_{1}}\varphi
_{j_{2}}\right\vert .
\end{align*}%
This gives the two particle conditions%
\begin{eqnarray*}
4\Gamma _{1}^{2}\left( \hat{\lambda}_{1}\right) ^{2} &=&2\Gamma
_{1}^{2}\left( \hat{\lambda}_{1}\right) \\
\left\{ \left( \lambda _{1j_{1}j_{1}}+\lambda _{1j_{2}j_{2}}\right) \right\}
^{2}-\left( \lambda _{1j_{1}j_{1}}+\lambda _{1j_{2}j_{2}}\right)
&=&0,\,1\leq j_{1}<j_{2}\leq 4
\end{eqnarray*}%
i.e. the matrix 
\begin{equation*}
\mathfrak{S}_{2}\left( \hat{\lambda}_{1}\right) ^{2}=\mathfrak{S}_{2}\left( 
\hat{\lambda}_{1}\right) , 
\end{equation*}%
where 
\begin{equation*}
\mathfrak{S}_{2}\left( \hat{\lambda}_{1}\right)
_{j_{1}j_{1}j_{2}j_{2}}=\lambda _{1j_{1}j_{1}}+\lambda
_{1j_{2}j_{2}},\,1\leq j_{1}<j_{2}\leq 4. 
\end{equation*}

If 
\begin{equation*}
\lambda _{1j_{1}j_{1}}=\lambda _{1j_{2}j_{2}}=1 
\end{equation*}%
then 
\begin{equation*}
\left\{ \left( 1+1\right) \right\} ^{2}-\left( 1+1\right) \neq 0, 
\end{equation*}%
while if 
\begin{eqnarray*}
\lambda _{1j_{1}j_{1}} &=&1 \\
\lambda _{1j_{2}j_{2}} &=&0
\end{eqnarray*}%
then 
\begin{equation*}
\left\{ 1\right\} ^{2}-\left( 1\right) =0,\,1\leq j_{1}<j_{2}\leq 4. 
\end{equation*}%
This shows that only one $\lambda _{1jj}$ can equal one while all other
values must zero thus $\hat{\lambda}_{1}$ is a one dimensional projector.

\section{Three Particle Condition}

The three particle expansion gives 
\begin{align*}
& \Gamma _{1}^{3}\left( \hat{\lambda}_{1}\right) =\hat{\lambda}_{1}\wedge P_{%
\mathcal{H}^{2}}=\sum_{\substack{ 1\leq j_{1}\leq 4  \\ 1\leq
j_{2}<j_{3}\leq 4}}\hat{\lambda}_{1j_{1}j_{1}}\left\vert \varphi
_{j_{1}}\right\rangle \left\langle \varphi _{j_{1}}\right\vert \wedge
\left\vert \varphi _{j_{2}}\varphi _{j_{3}}\right\rangle \left\langle
\varphi _{j_{2}}\varphi _{j_{3}}\right\vert \\
& =\frac{1}{3}\sum_{1\leq j_{1}<j_{2}<j_{3}\leq 4}\left( \lambda
_{1j_{1}j_{1}}+\lambda _{1j_{2}j_{2}}+\lambda _{1j_{3}j_{3}}\right)
\left\vert \varphi _{j_{1}}\varphi _{j_{2}}\varphi _{j_{3}}\right\rangle
\left\langle \varphi _{j_{1}}\varphi _{j_{2}}\varphi _{j_{3}}\right\vert
\end{align*}%
and its square is 
\begin{align*}
& \Gamma _{1}^{3}\left( \hat{\lambda}_{1}\right) ^{2}=\Gamma _{1}^{3}\left( 
\hat{\lambda}_{1}\right) \circ \Gamma _{1}^{3}\left( \hat{\lambda}%
_{1}\right) =\left( \hat{\lambda}_{1}\wedge P_{\mathcal{H}^{2}}\right) \circ
\left( \hat{\lambda}_{1}\wedge P_{\mathcal{H}^{2}}\right) \\
& =\frac{1}{9}\sum_{1\leq j_{1}<j_{2}<j_{3}\leq 4}\left\{ \left( \lambda
_{1}{}_{j_{1}j_{1}}^{2}+\lambda _{1}{}_{j_{2}j_{2}}^{2}+\lambda
_{1}{}_{j_{3}j_{3}}^{2}\right) \right. \\
& +\lambda _{1}{}_{j_{1}j_{1}}\lambda _{1}{}_{j_{2}j_{2}}+\lambda
_{1}{}_{j_{1}j_{1}}\lambda _{1}{}_{j_{3}j_{3}}+\lambda
_{1}{}_{j_{2}j_{2}}\lambda _{1}{}_{j_{1}j_{1}}+\lambda
_{1}{}_{j_{2}j_{2}}\lambda _{1}{}_{j_{3}j_{3}}+\left. \lambda
_{1}{}_{j_{3}j_{3}}\lambda _{1}{}_{j_{1}j_{1}}+\lambda
_{1}{}_{j_{3}j_{3}}\lambda _{1}{}_{j_{2}j_{2}}\right\} \left\vert \varphi
_{j_{1}}\varphi _{j_{2}}\varphi _{j_{3}}\right\rangle \left\langle \varphi
_{j_{1}}\varphi _{j_{2}}\varphi _{j_{3}}\right\vert
\end{align*}%
This gives the three particle conditions%
\begin{equation*}
\left\{ \left( \lambda _{1j_{1}j_{1}}+\lambda _{1j_{2}j_{2}}+\lambda
_{1}{}_{j_{3}j_{3}}\right) \right\} ^{2}-\left( \lambda
_{1j_{1}j_{1}}+\lambda _{1j_{2}j_{2}}+\lambda _{1}{}_{j_{3}j_{3}}\right)
=0,\,1\leq j_{1}<j_{2}<j_{3}\leq 4 
\end{equation*}%
i.e. the matrix 
\begin{equation*}
\mathfrak{S}_{3}\left( \hat{\lambda}_{1}\right) ^{2}=\mathfrak{S}_{3}\left( 
\hat{\lambda}_{1}\right) , 
\end{equation*}%
where 
\begin{equation*}
\mathfrak{S}_{3}\left( \hat{\lambda}_{1}\right)
_{j_{1}j_{1}j_{2}j_{2}j_{3}j_{3}}=\lambda _{1j_{1}j_{1}}+\lambda
_{1j_{2}j_{2}}+\lambda _{1}{}_{j_{3}j_{3}},\,1\leq j_{1}<j_{2}<j_{3}\leq 4. 
\end{equation*}%
It seems that 
\begin{equation*}
C_{3}\left( \hat{\lambda}_{1}\right) =\hat{\lambda}_{1}\wedge \hat{\lambda}%
_{1}\wedge \hat{\lambda}_{1}=\exp \left\{ \mathfrak{S}_{3}\left( \hat{\lambda%
}_{1}\right) \right\} 
\end{equation*}%
and in general 
\begin{equation*}
C_{N}\left( \hat{\lambda}_{1}\right) =\left( \hat{\lambda}_{1}\right)
^{\wedge N}=\exp \left\{ \mathfrak{S}_{N}\left( \hat{\lambda}_{1}\right)
\right\} , 
\end{equation*}%
where 
\begin{equation*}
\mathfrak{S}_{N}\left( \hat{\lambda}_{1}\right) _{j_{1}j_{1}\cdots
j_{N}j_{N}}=\sum_{1\leq \kappa \leq N}\lambda _{1j_{\kappa }j_{\kappa }}. 
\end{equation*}%
One also has%
\begin{equation*}
\mathfrak{S}_{N}\left( \hat{\lambda}_{1}\right) =\ln \left\{ C_{N}\left( 
\hat{\lambda}_{1}\right) \right\} 
\end{equation*}%
and the general condition 
\begin{equation*}
\mathfrak{S}_{N}\left( \hat{\lambda}_{1}\right) ^{2}=\mathfrak{S}_{N}\left( 
\hat{\lambda}_{1}\right) . 
\end{equation*}

\end{document}
